
\section{Definitions}

\subsection{Causality}

Causation is a subtle concept that cannot be fully
described using the language of Boolean logic [151] or that
of probabilistic inference; it requires the additional
notion of intervention [183], [237]. The manipulative
definition of causation [118], [183], [237] focuses on the
fact that conditional probabilities (“seeing people with
open umbrellas suggests that it is raining”) cannot reliably
predict the outcome of active intervention (“closing
umbrellas does not stop the rain”). Causal relations
can also be viewed as the components of reasoning
chains [151] that provide predictions for situations that
are very far from the observed distribution and may
even remain purely hypothetical [163], [183] or require
conscious deliberation [128]. In that sense, discovering
causal relations means acquiring robust knowledge that
holds beyond the support of observed data distribution
and a set of training tasks, and it extends to situations
involving forms of reasoning. 